\documentclass[11pt]{article}
\usepackage[margin=1in]{geometry}
\usepackage{amsmath,amssymb,amsthm}
\usepackage[hidelinks]{hyperref}

\newtheorem*{theorem}{Theorem}

\title{A Formulation Audit of Erdos Problem \#728}
\author{Codex sandbox note}
\date{July 7, 2026}

\begin{document}
\maketitle

\begin{abstract}
This note records a non-novel formulation audit of Erdos Problems \#728.
The literal written statement is true for a degenerate reason: it imposes
lower bounds on \(a,b\), but no upper bounds. The construction below should
not be read as a solution of the intended nontrivial problem.
\end{abstract}

\section*{Status}

This is an internal sandbox finding, not claimed as novel, and not worth
publishing as a standalone mathematical result. It is useful only as a
statement-normalization check.

The source page is
\[
\text{\url{https://www.erdosproblems.com/728}}.
\]

\section*{Written Statement}

Problem \#728 asks whether, for sufficiently small constants \(C>0\) and
\(\epsilon>0\), there are infinitely many integers \(a,b,n\) such that
\[
a\geq \epsilon n,\qquad b\geq \epsilon n,\qquad
a!b!\mid n!(a+b-n)!,
\]
and
\[
a+b>n+C\log n.
\]

\section*{Result}

\begin{theorem}
For every fixed \(C>0\) and every fixed \(\epsilon>0\), the written
statement of Problem \#728 has infinitely many solutions.
\end{theorem}

\begin{proof}
Fix \(C>0\) and \(\epsilon>0\). For any sufficiently large positive integer
\(n\), choose an integer
\[
w\geq \max(C\log n,\epsilon n).
\]
Define
\[
a=n+w+1,\qquad b=\frac{(n+w+1)!}{n!}-1.
\]
Then \(a\geq \epsilon n\), and \(b\geq \epsilon n\) for all sufficiently large
\(n\). Also
\[
a+b-n=b+w+1.
\]
Now compute
\[
\frac{n!(a+b-n)!}{a!b!}
=
\frac{n!(b+w+1)!}{(n+w+1)!b!}.
\]
Since
\[
\frac{(n+w+1)!}{n!}=b+1,
\]
we obtain
\[
\frac{n!(a+b-n)!}{a!b!}
=
\frac{(b+1)(b+2)\cdots(b+w+1)}{b+1}
=(b+2)(b+3)\cdots(b+w+1),
\]
which is an integer. Hence \(a!b!\mid n!(a+b-n)!\).

Finally,
\[
a+b-n=b+w+1>C\log n,
\]
so \(a+b>n+C\log n\). Since \(n\) can be chosen arbitrarily large, there
are infinitely many such triples.
\end{proof}

\section*{Interpretation}

The construction uses \(a>n\) and huge \(b\). Therefore it does not touch
the intended logarithmic-barrier problem. It only shows that the written
statement is too weak unless additional upper bounds are included, for
example \(a,b\leq (1-\delta)n\).

Even adding only \(a,b\leq n\) would not fully remove degeneracy, since
\(a=b=n\) gives \(n!n!\mid n!n!\) and \(2n>n+C\log n\) for all sufficiently
large \(n\).

\section*{Verdict}

This is a useful formulation audit, but not a publishable new finding.
The active search for publishable work should move to a genuinely open
target or to a strengthened version of \#728 with the intended hypotheses.

\end{document}
